\section{Machine verification}
\label{sec:lean}

A substantial part of this paper is formalised in Lean~4 and checked by its kernel. This
section reports exactly what is and is not machine-verified. Toolchain:
\texttt{leanprover/lean4:v4.33.0-rc1}, Mathlib revision
\texttt{cd580e54f1a6b46063824e80cec92f64692cbe78}.

\subsection{What is formalised}

\begin{center}\small
\begin{tabular}{@{}ll@{}}
\toprule
result & status\\
\midrule
Theorem~\ref{thm:arow} (first row, Ap\'ery) & formalised, $0$ sorries\\
first row for the Brown--Zudilin double sum $Q_n$ & formalised, $0$ sorries\\
full digit-product form of both first rows & formalised, $0$ sorries\\
Lemma~\ref{lem:K} and Corollary~\ref{cor:K2} & formalised, $0$ sorries\\
Theorem~\ref{thm:LB} (general $\chi$-twisted descent) & formalised, $0$ sorries\\
%% FIXED [MINOR]: the Lean theorem is about $B_{\min}$, not about $b$; Theorem \ref{thm:brow}
%% FIXED is the statement for $b$, and the bridge $B_{\min}=b$ is the unformalised row below.
second row for the minimal form $B_{\min}$ (the $B_{\min}$ analogue of
Theorem~\ref{thm:brow}) & formalised, $0$ sorries\\
second row for Franel & formalised, $0$ sorries\\
two- and three-digit Kummer module (K1--K4) & formalised, $0$ sorries\\
\midrule
$B_{\min}=b$ (Theorem~\ref{thm:minimal}) & \emph{not} formalised (proved on paper)\\
Conjecture~\ref{conj:master} (master form) & not formalised (not proved)\\
\bottomrule
\end{tabular}
\end{center}

\texttt{lake build} is clean; a clean rebuild of all eight modules takes $34$ seconds with
Mathlib prebuilt. A grep for \texttt{sorry}, \texttt{admit} and \texttt{native\_decide}
over the development returns nothing, and \texttt{\#print axioms} on every exported
theorem returns only the three standard Mathlib axioms:
\begin{verbatim}
'ZetaLucas.K_descent'     depends on axioms:
    [propext, Classical.choice, Quot.sound]
'ZetaLucas.mon_descent'   depends on axioms: [ ... same three ... ]
'ZetaLucas.theorem_LB'    depends on axioms: [ ... same three ... ]
'ZetaLucas.bMin_lucas'    depends on axioms: [ ... same three ... ]
'ZetaLucas.bFranel_lucas' depends on axioms: [ ... same three ... ]
'ZetaLucas.apery_lucas'   depends on axioms: [ ... same three ... ]
'ZetaLucas.Q_lucas'       depends on axioms: [ ... same three ... ]
\end{verbatim}
In particular there is no \texttt{sorryAx} and no \texttt{Lean.ofReduceBool}: no
statement rests on a compiled-code evaluation.

\subsection{The statements}

We display the main statements in ASCII-transliterated Lean (Unicode identifiers such as
$\mathbb Q$, $\mathbb N$, $\in$, $\forall$ are written \texttt{Q}, \texttt{N},
\texttt{in}, \texttt{forall}). The predicate \texttt{PCong p x y} is defined as
\texttt{padicNorm p (x - y) < 1}, i.e.\ $x\equiv y\pmod p$ inside $\Zp\subseteq\Q$; it
implies $p$-integrality of both sides as soon as one of them is $p$-integral, so the
statements are not vacuous.

\begin{verbatim}
-- Lemma K
theorem K_descent {chi : N -> Z} (hchi : forall m n, chi (m*n) = chi m * chi n)
    {r : N} (hr : 0 < r) (y : N) :
    PCong p ((p : Q)^r * K chi r y) ((chi p : Q) * K chi r (y / p))

-- Theorem LB (hypotheses H1, H2, H3, H4, H4c, Hw, H5 as in Section 4)
theorem theorem_LB {J : Finset iota} {c : iota -> Q}
    {mon : iota -> List Letter} {S : N -> N -> Z} {w : N} {chip : Z}
    (H1 ...) (H2 ...) (H3 ...) (H4 ...) (H4c ...) (Hw ...) (H5 ...)
    {a r : N} (ha : a < p) (hr : r < p) :
    PCong p ((p : Q)^w * Brow J c mon S (a*p + r))
      ((chip : Q) * Brow J c mon S a * (Arow S r : Q))

-- the second row for the minimal Apery form; bMin n is the sum over
-- k in range (n+1) of (A n k : Q) * (2 * Harm 3 n - Harm 3 k)
theorem bMin_lucas {p : N} [Fact p.Prime] {a r : N}
    (ha : a < p) (hr : r < p) :
    PCong p ((p : Q)^3 * bMin (a*p + r)) (bMin a * (apery r : Q))

-- the second row for Franel
theorem bFranel_lucas {p : N} [Fact p.Prime] (hp2 : p <> 2) {a r : N}
    (ha : a < p) (hr : r < p) :
    PCong p ((p : Q)^2 * bFranel (a*p + r)) (bFranel a * (franel r : Q))
\end{verbatim}

\subsection{Two sharpenings found by formalising}

\begin{enumerate}[leftmargin=1.7em]
\item \texttt{bMin\_lucas} carries \emph{no} hypothesis $p\ge5$ and \emph{no} hypothesis
  $1\le a$. The minimal weight's coefficients are the integers $2,-1$, so the (H4)
  coefficient condition is vacuous and the theorem holds for every prime, including $2$
  and $3$; and $a=0$ is fine, since $B_{\min}(0)=0$ and $p^3B_{\min}(r)$ is $p$-divisible
  for $r<p$. The informal Theorem~\ref{thm:brow} needs $p\ge5$ only because Ap\'ery's
  original weight carries $1/(2m^3)$.
\item \texttt{bFranel\_lucas} needs only $p\ne2$, where the informal analysis had recorded
  $p\ge5$; the coefficient denominators are $4$, and $p=3$ is fine.
\end{enumerate}

\subsection{Design decisions, and why they matter mathematically}

Three changes were made to the informal hypotheses; all three make the theorem easier to
instantiate, not harder.

\begin{enumerate}[leftmargin=1.7em]
\item (H1) is stated as an unconditional congruence
  $S(ap+r,bp+s)\equiv S(a,b)S(r,s)\pmod p$ for \emph{all} $a,b$, together with
  $S(n,k)=0$ for $k>n$; (H2) then disappears. In the carrying regime both sides vanish
  modulo $p$, so the dichotomy is absorbed. The consequence is that \emph{no
  product-region argument occurs anywhere in the formalisation}: the double sum factors
  because a full block $\{0\le k<(a+1)p\}$ splits, and the extension from
  $\{0\le k\le n\}$ to that block is an equality of integers. This retires, for the general
  theorem, the compression that Lemma~\ref{lem:fibre} had to expand by hand.
\item (H3)'s ``surviving set'' is replaced by the side condition
  $S(r,s)\not\equiv0\pmod p$, which is exactly what excludes a borrow in an argument such
  as $n-k$: for Franel it gives $\bin rs\ne0$, hence $s\le r$, hence
  $\lfloor(n-k)/p\rfloor=a-b$. Instances then \emph{derive} the no-borrow condition rather
  than describing $\Sigma_r$.
\item (H5) is stated as ``the product of the letters' characters at $p$ is the same for
  every monomial'' rather than as a count of $\chi$-letters. This is strictly more
  general --- it allows mixed conductors --- and cheaper to check.
\end{enumerate}

The letter system is deliberately concrete: a \texttt{Letter} bundles a completely
multiplicative $\chi:\N\to\Z$, a positive degree, and an argument form
$\N\to\N\to\N$; a monomial is a \texttt{List Letter}, so genuine products such as
$H_kH_{n-k}$ (which Franel needs) are expressible. The sum $\Kchi ry$ is taken over
\texttt{range (y+1)}, so the $m=0$ term is $\chi(0)/0^r=0$ in $\Q$ for $r\ge1$; this makes
the reindexing of the $p$-divisible layer a single bijection $m\mapsto m/p$,
$j\mapsto jp$, and no reindexing of $\texttt{Finset.Icc}$ --- the step that had been
priced as the main schedule risk --- ever arises.

\subsection{Numerical checks inside the build}

Normalisations are pinned in-file, by kernel evaluation where feasible and by
\texttt{\#eval} otherwise:
$\texttt{apery}\mapsto1,5,73,1445,33001,819005$ (A005259, by \texttt{decide});
$Q\mapsto1,21,2989,714549,217515501,76157194521$, confirming the double-sum reading
$Q_1=21$; $\texttt{bMin}\mapsto0,6,\tfrac{351}4,\tfrac{62531}{36},\tfrac{11424695}{288},
\tfrac{35441662103}{36000}$, which is exactly the classical sequence $b_n$;
$\texttt{franel}\mapsto1,2,10,56,346,2252$ (A000172) and
$\texttt{bFranel}\mapsto0,1,4,\tfrac{208}9,\tfrac{1280}9$, matching the Franel recurrence
with $B(0)=0$, $B(1)=1$. An evaluable form of $p$-divisibility sweeps every cell:
both second-row congruences evaluate to \texttt{true} for all $a,r<p$ at $p=5,7$ in the
build, and off-line at $p=11,13$. Sharpness spot-check: at $p=5$, $a=r=1$ the difference
$p^3B_{\min}(ap+r)-B_{\min}(a)\,a_r$ equals $103188530395/32$, of $5$-valuation exactly
$1$ --- the floor is exactly $w$, as the sweeps of Section~\ref{sec:depth} report.

\subsection{Effort, and the economics of the two routes}

The development is $1540$ lines across eight modules
(\texttt{Core}, \texttt{Apery}, \texttt{BrownZudilin}, \texttt{PadicBridge},
\texttt{Letters}, \texttt{TheoremLB}, \texttt{Instances}, \texttt{Kummer}).

The route matters more than the rate. Formalising Theorem~\ref{thm:brow} \emph{via the
proof of Section~\ref{sec:weight3}} --- Lemma~\ref{lem:V}, Corollary~\ref{cor:tailfact},
Lemma~\ref{lem:levelcarry}, the Kummer ledger and the block reindexing of
$\texttt{Finset.Icc}$ --- was costed in advance at $\approx13$ hours $\pm4$, with two
pieces carrying essentially all the schedule risk. Formalising it \emph{via
Theorem~\ref{thm:LB} and the minimal form} cost $\approx2.2$ hours and produced a strictly
stronger result: arbitrary $\chi$, arbitrary weight $w$, all primes, $a=0$ allowed, and a
second sequence for free. The claim made informally --- that the general theorem
``reproves the weight-three theorem in three lines, with no Lemma~V, no tail
factorisation and no Kummer ledger'' --- is thus machine-confirmed, and the cost ratio is
about $6\times$. (Wall-clock figures are reconstructed from the work order rather than
stopwatch readings and should be read as $\pm25\%$; line counts and build-iteration counts
are exact.)

The Kummer module is retained even though the $(\mathrm{LB}_w^\chi)$ route does not use
it: it is precisely the toolkit that the non-tame instances of
Section~\ref{ssec:lemmaD} will need.

\subsection{What is \emph{not} formalised}

\begin{enumerate}[leftmargin=1.7em]
\item \emph{$B_{\min}=b$}. The Lean theorem is about the explicit sum
  $B_{\min}(n)=\sum_kA(n,k)(2\HH3n-\HH3k)$. Its equality with the classical Ap\'ery
  numerator is Theorem~\ref{thm:minimal}, proved on paper here and \emph{not} formalised.
  The gap is therefore a formalisation gap, not a mathematical one, and it is well-scoped:
  Lemma~\ref{lem:conv}'s six conversions, the six certificates of Table~\ref{tab:certs}
  (each an \texttt{Expand}$\to0$ polynomial identity, i.e.\ \texttt{ring} in Lean) and an
  induction on $L$, whose leading coefficient $(n+1)^3$ never vanishes.
\item \emph{$B_{\mathbf A}$ is the second Franel solution}: same status --- it means
  proving that the explicit sum satisfies the Franel recurrence, which is
  Theorem~\ref{thm:franel} here.
\item \emph{The master form}, Conjecture~\ref{conj:master}: not proved, hence not
  formalised. Note that the first-row multi-digit form \emph{is} formalised and is a cheap
  induction; the second row is not, because its ratio form has poles.
\item \emph{A twisted instance.} The $\chi$ machinery of \texttt{theorem\_LB} is proved
  but never exercised at $\chi(p)\ne1$, because the only known twisted sequence
  ($\mathbf E$) is non-tame (Problem~\ref{prob:lemmaD}).
\item \emph{The non-tame instances} and \emph{the seven families with no known
  decomposition}: nothing to formalise until the mathematics exists.
\end{enumerate}
